Write \(b_t=B_{\cdot,t}\), \(h_t=h_e(t)\), and form the finite observable--target hull
\[
\mathcal C=\left\{(Bv,h_e^{\top}v):v\in\Delta_{\mathrm{SLCC}}\right\}
=\operatorname{conv}\{(b_t,h_t):t\in\mathcal T_{\mathrm{SLCC}}\}.
\tag{1}
\]
The simplex is compact and the maps in (1) are linear.  Hence \(\mathcal C\), every fiber \(\{v\in\Delta_{\mathrm{SLCC}}:Bv=p\}\), and its image under \(v\mapsto h_e^{\top}v\) are compact.  For \(p\in\mathfrak O\), the last image is a nonempty compact interval, because the fiber is convex.  Its endpoints are therefore attained and equal \(L(p)\) and \(U(p)\) by \(\text{def:sharp-endpoint-programs}\).

We derive a finite half-space representation rather than assuming a normalization of a dual variable.  The membership condition \((x,z)\in\mathcal C\) is the finite linear system
\[
v_t\geq0\ (t\in\mathcal T_{\mathrm{SLCC}}),\qquad
\mathbf 1^{\top}v=1,\qquad Bv=x,\qquad h_e^{\top}v=z.
\tag{2}
\]
Eliminate the coordinates of \(v\) successively.  At one elimination step, inequalities with positive coefficient of the eliminated coordinate give upper bounds, inequalities with negative coefficient give lower bounds, and inequalities with zero coefficient are retained; existence of that coordinate is equivalent to every retained lower bound being no larger than every retained upper bound.  Thus one step replaces a finite system by a finite system and is an equivalence.  Induction over the finitely many coordinates in (2) gives finite sets \(R\), vectors \(a_r\), and scalars \(d_r,c_r\) such that
\[
\mathcal C=\{(x,z):a_r^{\top}x+d_rz\leq c_r\text{ for every }r\in R\}.
\tag{3}
\]
The valid inequalities \(-z\leq0\) and \(z\leq1\), which follow from \(h_t\in\{0,1\}\), may be appended to (3).  Hence the index sets \(R_-:=\{r:d_r<0\}\) and \(R_+:=\{r:d_r>0\}\) are nonempty.

For fixed \(p\in\mathfrak O\), the inequalities with \(d_r=0\) are already satisfied, and (3) says exactly
\[
L(p)=\max_{r\in R_-}\frac{c_r-a_r^{\top}p}{d_r},
\qquad
U(p)=\min_{r\in R_+}\frac{c_r-a_r^{\top}p}{d_r}.
\tag{4}
\]
For \(r\in R_-\), write the affine function in (4) as \(\alpha_r+\lambda_r^{\top}p\).  Substitution of each generator \((b_t,h_t)\) into its defining inequality gives
\[
\alpha_r+\lambda_r^{\top}b_t\leq h_t\quad\text{for every }t,
\tag{5}
\]
so \((\alpha_r,\lambda_r)\in\mathcal D_L\).  Conversely, if \((\alpha,\lambda)\in\mathcal D_L\) and \(Bv=p\), \(\mathbf 1^{\top}v=1\), then
\[
\alpha+\lambda^{\top}p
=\sum_t v_t(\alpha+\lambda^{\top}b_t)
\leq\sum_t v_th_t=h_e^{\top}v.
\tag{6}
\]
Taking the minimum in (6), then choosing an index attaining the finite maximum in (4), proves
\[
L(p)=\max_{(\alpha,\lambda)\in\mathcal D_L}
(\alpha+\lambda^{\top}p),
\tag{7}
\]
with a nonempty optimal face.  The same argument with \(r\in R_+\) reverses (5)--(6) and proves
\[
U(p)=\min_{(\alpha,\lambda)\in\mathcal D_U}
(\alpha+\lambda^{\top}p),
\tag{8}
\]
again with attainment.

Equation (4) represents \(L\) as the maximum of finitely many affine functions and \(U\) as the minimum of finitely many affine functions.  They are therefore finite, piecewise affine, and continuous at every point of \(\mathfrak O\), including its relative boundary.  Repeating (1)--(8) with the finite matrix \(K_{\mathrm{blind}}B\) proves the corresponding assertions for \(L_{\mathrm{blind}}\) and \(U_{\mathrm{blind}}\) on \(K_{\mathrm{blind}}\mathfrak O\).